\subsection*{Biological datasets and support--query separation}

PBMC evaluation used the 3-prime RNA component of the Hao et al. CITE-seq reference\cite{hao2021multimodal}: GEO accession GSE164378, RNA sample GSM5008737 and the supplied cell metadata. The source contains 161,764 cells and 33,538 RNA coordinates. Published annotations combine RNA, measured surface proteins and curation; they are not an independent protein-only ground truth. Predictions used RNA alone. Barcodes joined the original RNA matrix to metadata, and RNA library sums were checked against the metadata. The accompanying antibody measurements and annotations were audited against the author-linked reproduction snapshot; antibody measurements were not prediction inputs. Cells annotated as doublets were excluded. A class was eligible only when it contained at least 25 cells in each support donor and ten in each query donor, yielding eight broad and twenty finer classes.

Support donors were P1, P3, P5 and P7; query donors were P2, P4, P6 and P8. All timepoints of a donor remained on the same side, with two donors from each acquisition batch on each side. Ten query cells per class and donor were selected by deterministic metadata-based SHA256 rankings, giving 320 and 800 query cells for the broad and fine tasks. Within each class and support draw, donor and cell rankings were shuffled deterministically and sampled round-robin across donors. Prefixes supplied one, ten or 100 support cells per class; at 100 supports, each donor contributed 25 cells. Query cells were fixed across support budgets, draws and methods. The primary comparison used support draws 0--2; extended diagnostics followed the same sampling rule with more draws.

Human cortex evaluation used the 50,000-cell ZEBRA benchmark subsample\cite{flotho2024zebra}, with 21,158 genes and four broad classes: glia, immune, neural and vascular cells. The shared support pool comprised Bakken, Gerrits, Nagy, Schirmer, Trobisch and Yang Cortex studies. Development queries comprised 7,769 cells from 59 recorded donors in Absinta, Kihara, Lim and Velmeshev studies; held-out-study queries comprised 7,163 cells from 34 recorded donors in Hardwick, Herring, Jaekel and Yang studies. Support and query studies were disjoint, and query cells remained fixed across budgets and methods. Both query panels nevertheless contributed to model development; the held-out-study designation describes their study separation, not an untouched final model-selection set. Frozen cell identifiers, study assignments and input-file hashes define these panels more precisely than publication names alone.

The secondary kidney evaluation used healthy living-donor data from GSE202109\cite{mcevoy2022kidney}, distributed through CELLxGENE, with 27,675 cells from 19 donors and 26,791 measured genes. Three donor-held-out folds evaluated five- and nine-class cell-type tasks and proximal-tubule sex prediction. Cell-type tasks supplied 100 supports per class and 40 query cells per class. Sex prediction used six or seven query donors per fold, with five query cells per donor. Two 4,096-gene panels were retained: a kidney-support-derived panel and a panel derived from the historical ZEBRA support pool. These feature panels therefore used more support information than an individual few-shot case. Kidney cell-type accuracy and donor-level sex accuracy were kept separate.

\subsection*{Expression inputs and feature selection}

CellPFN received original counts for PBMC and the supplied processed expression for the primary cortex analysis. Count-based comparator inputs were transformed to $\log[1+10^4x_{ig}/\sum_hx_{ih}]$, using each cell's full measured library before gene selection. Supplied ZEBRA log-transformed counts per 10,000 were passed through unchanged. Genes unavailable in individual ZEBRA studies had been stored as zeros; this preprocessing artifact was evaluated separately from biological non-expression.

Highly variable genes were selected using labelled support rows only. The implementation exponentiated log-expression, computed gene means and unbiased variances, and standardized log dispersion within 20 mean-expression bins, following the Seurat-flavour dispersion statistic. Finite scores were ranked with a deterministic tie break based on the anonymous support column, retaining at most 2,000 or 8,192 genes. The same selected columns and order were applied to queries. Table labels denote requested HVG caps. In the three primary one-shot draws, the 8,192-HVG view retained 6,995--7,647 genes for broad PBMC and 2,209--5,301 for cortex; finite support variation could therefore reduce the realized width below the requested cap. Principal-component controls standardized support-selected genes, clipped standardized values to $[-10,10]$, fitted the decomposition to support rows, and retained $\min(100,n_{\mathrm{support}}-1,n_{\mathrm{genes}})$ components. Query expression did not influence selection, scaling or projection.

The cortex sensitivity analysis reconstructed integer counts from the lattice of the supplied processed values and retained 12,325 genes measured across eligible contributing studies. Reconstruction was conditional on the smallest positive expression level representing one count, rather than a retrieval of original sequencing counts. A matched processed partner was recomputed from the reconstructed common panel. Query identities and labels were unchanged; some supports changed after reconstruction-ineligible cells were removed. Consequently, historical full-panel and common-panel results do not isolate count versus processed representation. The reconstruction rule and support replacement are specified in the technical supplement.

\subsection*{Comparator models and selection}

Released tabular comparators were TabICLv2, TabPFN-3.5 and Kumo Tabular small and large\cite{qu2026tabiclv2,jaeger2026tabpfn35,nvidia2026kumo}. Each used labelled support cells as context, with no target-task gradient update. Continued-training checkpoints used the same inference route as their released reference. Main comparisons used support-selected 2,000-gene inputs; scaling analyses separately evaluated 8,192-gene and all-gene inputs. Package-default transformations, feature subsampling and feature limits remained active. Refused inputs and unevaluated cases were recorded as missing, not as zero accuracy. TabPFN-3.5 results were obtained under its non-commercial research terms\cite{priorlabs2026tabpfn35card}.

Geneformer-V2-316M\cite{chen2026geneformer} and the whole-human scGPT checkpoint\cite{cui2024scgpt} supplied frozen embeddings. Geneformer used Ensembl-linked rank-value tokens and mean embeddings; scGPT used gene-symbol inputs, binned values and unit-length CLS embeddings. Logistic classifiers were fitted to support embeddings only. Geneformer used support-standardized embeddings; scGPT used unscaled embeddings. Exact token limits, classifier settings and recorded software releases are listed in Table~\ref{tab:evaluation-settings}. Native input routes were preserved, so these comparisons jointly reflect representation, feature coverage and classifier choice.

Support-trained comparators included nearest-class-mean and prototype rules, logistic regression, diagonal and low-rank LDA, CellTypist\cite{dominguezconde2022celltypist}, SingleR\cite{aran2019singler} and scANVI\cite{xu2021scanvi}. CellTypist was fitted to the selected supports without query majority voting. SingleR used the supports as its reference. Inductive scANVI used supports only; its separately labelled transductive comparison also used unlabelled query cells and a donor batch key. Processed cortex inputs were not passed to count-likelihood scANVI. LDA diagnostic rows were available at ten and 100 supports, not one support per class. The inductive scANVI recipe used batch size 128, a 90\% training fraction and ten supervised epochs without class-balanced replacement sampling. At fine-PBMC one-shot support, 18 of 20 labelled cells entered optimization, leaving at least two classes without a training example. These low-support results describe that fixed recipe. SingleR retained its default classic marker mode rather than the alternative recommended for single-cell references.

The displayed family representatives were chosen within the development evaluation. Kumo size selection used the mean across the twelve core panel--budget combinations on 2,000 genes; the CellPFN reference was an explicitly selected model state, and logistic heads were the declared primary heads for Geneformer and scGPT. These choices were not independently validated model-selection procedures. Study-level pretraining overlap was assessed separately from benchmark support--query separation, using the recorded corpus/source-study audit. Such matches indicate possible unlabelled query-source exposure; they do not establish identical cell inclusion, target-label exposure or supervised leakage.

\subsection*{Exploratory subtype and sex tasks}

The eight PBMC subtype tasks comprised level-2 CD4 and combined T-cell siblings, level-3 CD4, CD8, B, NK and combined T-cell siblings, and the complete eligible level-3 vocabulary. They inherited the PBMC donor split and class-availability rules. Candidate selection was conducted after the prior-ablation read-outs and is exploratory. It used nearest-class-mean and logistic-reference performance, not CellPFN or proxy predictions: at 100 supports, the larger reference accuracy had to reach a chance-corrected skill of 0.20 while remaining below 85\% balanced accuracy. Each task used ten support draws at one and ten supports and five draws at 100 supports. CellPFN received up to 8,192 support-selected genes; the Kumo-small proxy received up to 2,000. Supplied CellPFN task means were reported without constructing intervals from unavailable individual-case predictions. Two exploratory cortex tasks used the same reference checkpoint and up to 8,192 HVGs: seven development labels and three held-study neural labels, defined in Supplementary Table~\ref{tab:s9}. They used ten, ten and five support draws at one, ten and 100 supports; the supplied checkpoint means do not include draw-level uncertainty. Changes in labels and evaluated cells preclude interpreting their difference from the primary broad-class benchmark as an isolated difficulty effect.

The original ZEBRA sex task used Sayed oligodendrocytes, 100 support cells per sex from twelve GSM samples, and 50 query cells from ten other GSM samples. GSM separation was not interpreted as independently verified donor separation. Later study-matched sex tasks used ten splits, four to 32 support units per sex, ten cells per support unit and twenty per query unit. The positive control required measurement of XIST and at least four specified Y-linked genes; a broader task included additional studies with a different common gene universe. Units were persons except for age-by-sex groups of GSM samples in Sayed. Reused persons and unresolved cross-study identity overlap preclude treating the ten splits as independent cohorts. Scores were averaged within query units before unit-level AUC calculation. \unresolvedclaim{O06}{Gene-level attribution for these sex tasks remains unavailable; eligibility based on sex-linked genes does not demonstrate that those genes drove the model's predictions.}

\subsection*{Metrics and uncertainty}

Balanced accuracy was the mean recall across every declared query class. Negative log-likelihood was the mean negative log probability of the true class, in nats; expected calibration error used fifteen fixed-width confidence bins. Probabilities and class order were checked before scoring. Performance means first averaged support draws within a panel--budget combination and then weighted the twelve combinations equally. Low-support summaries averaged the eight combinations with one or ten supports. Paired comparisons matched model, panel, input view, support budget and draw; a difference between unmatched views was not interpreted as an isolated model effect.

Three-draw benchmarks are descriptive and use ranges or individual values. The ten- and twenty-draw diagnostics used Student $t$ intervals over support draws, paired for contrasts. These quantify sensitivity to labelled-context sampling conditional on the chosen cohorts and weights. Every trained arm had one model seed; task-stream replicates did not estimate variation across independently initialized models. Kidney fold differences combine cohort and support variation. The nested PBMC subtype tasks share cells and labels and are not independent biological replicates. Attribution case ranges and single-run timings were not assigned confidence intervals. Missing, refused and unrun values were never replaced by zero.

\subsection*{Query-gene deletion and attribution validation}

For a query $x$ and fixed support set $S$, let $a$ and $b$ be the initially highest- and second-highest-logit classes. Define the fixed-pair margin $m(x;S)=f_a(x;S)-f_b(x;S)$. The exact deletion target for gene $g$ was
\begin{equation}
 I_g(x;S)=m(x;S)-m(x^{(g\leftarrow0)};S).
\end{equation}
Only the query gene was set to zero; support summaries and the fitted covariance operator were held fixed. A positive effect supports the original predicted class relative to its original runner-up. It is neither an expression-direction statistic nor an experimental perturbation effect. Negative effects oppose this fixed-pair margin. A gene already at zero has zero deletion effect, so this intervention cannot measure predictive information carried by its absence.

DS-LOO updates the mean, root-mean-square, maximum and log-sum-exp pools when each per-gene relation vector is replaced by its zero-expression counterpart. A first-order expansion through the prediction head maps these pooled changes to margin changes. An additional gradient calculation corrects the read-depth coupling in count mode and the detected-gene fraction. This is an attribution sweep with head differentiation and corrections, rather than one ordinary prediction or exact model evaluation per gene. The supported mathematical approximation is detailed in the technical supplement. \unresolvedclaim{M01}{The production attribution implementation, exact normalization corrections and treatment of query-side operator coupling remain to be reconciled with the recorded analysis before the procedure is fully reproducible.}

Validation used six case configurations and both the reference and no-operator model: three PBMC cases, one cortex case, and one kidney case evaluated in count and processed modes. Cells had to be correctly classified under both stored bfloat16 predictions and the float32 rerun. Seeded sampling selected at most 24 cells per class, or 40 for kidney. Exact deletion was evaluated on unions of method-selected genes and random expressed genes, with additional DS-LOO extreme-ranked genes; it was not exhaustive over all genes. Rank correlations were computed within cells, then summarized by the recorded class/case aggregation; relative Euclidean error compared estimated and exact effect vectors. Deletion curves recomputed the fixed-pair margin after removing sets of ranked query genes. Gene symbols were restored after prediction for descriptive marker comparisons. No marker enrichment test was performed. Timings distinguished measured attribution runs from extrapolated exhaustive-deletion costs; comparisons made under shared GPU load remain provisional.

\subsection*{Provisional neuropathology analysis}

\unresolvedclaim{O02}{The ROSMAP analysis used aggregate records for donor-disjoint splits of neuropathology contrasts\cite{bennett2018rosmap}: CERAD 4 versus 1, CERAD 3--4 versus 1--2, and Braak 0--II versus V--VI. Support and query allocation was matched within sex on source and age bin. Each support donor contributed eighteen cells and each query donor twenty. Donor-level AUC was summarized over ten overlapping random splits using the reported corrected repeated-split interval. The exact sequencing-release crosswalk, donor-label corrections and CellPFN input-mode record remain unresolved, and these analyses must be rerun after cohort definitions are frozen.}
